% Contenido del póster: Larios-Hernandez (2025), Journal of Small Business & Entrepreneurship. DOI 10.1080/08276331.2024.2389505
% Elaborado con apoyo de IA a partir del resumen oficial del artículo (OpenAlex), 2026-11-12.
\def\DOI{10.1080/08276331.2024.2389505}
\def\TituloEs{Facilitar nuevos servicios financieros en América Latina: un análisis exploratorio del modelo de espacio de prueba regulatorio}
\def\TituloOriginal{Facilitating new financial services in Latin America: An exploratory analysis of the sandbox model}
\def\Autores{Larios-Hernandez}
\def\Anio{2025}
\def\Revista{Journal of Small Business \& Entrepreneurship}
\def\TituloGrande{\fontsize{58}{70}\selectfont}  % título largo: letra algo menor para que el póster quepa en una página

\def\Problema{Los espacios de prueba regulatorios (\textit{sandboxes}) permiten a las fintech probar productos bajo supervisión, pero la investigación los ha estudiado sobre todo desde el procedimiento legal y no desde la gestión.}
\def\Objetivo{Analizar de principio a fin los espacios de prueba para fintech, con énfasis en América Latina, desde una perspectiva de gestión, y proponer un marco teórico aplicable.}
\def\Metodo{Primero se caracterizan los motivos, las condiciones y los resultados esperados de estos espacios. Luego se comparan casos, a partir de recursos web, con la literatura. Por último, se aplica un método de Gioia invertido para caracterizar los modelos de espacio de prueba.}
\def\Resultados{\begin{itemize}
  \item América Latina sigue un enfoque fragmentado y gradual.
  \item Los espacios de prueba buscan eficiencias en pagos, inversión y crédito.
  \item Predomina la intermediación de flujos de capital.
  \item Participan tanto empresas establecidas como emprendimientos nuevos.
\end{itemize}}
\def\Conclusion{El diseño del espacio de prueba influye en el desarrollo de las fintech. El marco propuesto sirve a emprendedores y reguladores, y puede aplicarse a otras regiones.}
\def\Relevancia{El financiamiento participativo financiero es un caso de intermediación de flujos de capital, el tipo de actividad que, según el estudio, predomina en la experimentación regulatoria latinoamericana. En el Perú se rige por un régimen propio de la SMV desde el Decreto de Urgencia N.\textsuperscript{o}~013-2020, y el estudio ofrece criterios para evaluar cómo el diseño regulatorio favorece o frena la entrada de nuevas plataformas.}
\def\PalabrasClave{espacio de prueba regulatorio, fintech, América Latina, regulación, emprendimiento}
% --- Ficha del estudio (herramientas/posteres/ficha_articulos.py) ---
\def\Ambito{América Latina}
\def\Enfoque{Cualitativo exploratorio (método de Gioia invertido)}
\def\FiguraPeru{1}
% --- fin de la ficha ---
